% Source: https://artofproblemsolving.com/community/c3h1974978_the_problem
% Stable identity: aops-topic-1974978-q01
\textbf{Step 1: $\dfrac{ab}{a+2b}\le\dfrac{2a+b}{9}$.} Multiplying out, this is $9ab\le(2a+b)(a+2b)=2a^2+5ab+2b^2$, i.e. $2(a-b)^2\ge0$. Adding the three cyclic versions,
\[
P\le\dfrac{(2a+b)+(2b+c)+(2c+a)}{9}=\dfrac{a+b+c}{3}.
\]
\textbf{Step 2: $a+b+c\le3$.} Each of $a^2,b^2$ is less than $4$. Viewing the condition as a quadratic in $c$, $c^2+ab\,c+(a^2+b^2-4)=0$, its positive root is
\[
c=\dfrac{-ab+\sqrt{a^2b^2-4a^2-4b^2+16}}{2}=\dfrac{-ab+\sqrt{(4-a^2)(4-b^2)}}{2}.
\]
By AM--GM, $\sqrt{(4-a^2)(4-b^2)}\le\dfrac{8-a^2-b^2}{2}$, so
\[
c\le\dfrac{8-a^2-b^2-2ab}{4}=2-\dfrac{(a+b)^2}{4}.
\]
With $s=a+b$ this gives $a+b+c\le s+2-\dfrac{s^2}{4}=3-\dfrac{(s-2)^2}{4}\le3$.

Hence $P\le1$. For $a=b=c=1$ the condition holds ($1+1+1+1=4$) and $P=3\cdot\dfrac{1}{3}=1$. So the maximum is $\boxed{1}$.

\textbf{No minimum.} Always $P>0$. Take $b=c=t$ small and $a$ the positive number with $a^2+2t^2+at^2=4$; as $t\to0^+$, $a\to2$ and
\[
P=\dfrac{at}{a+2t}+\dfrac{t^2}{3t}+\dfrac{ta}{t+2a}\to0 .
\]
So $P$ takes values arbitrarily close to $0$ but never $0$: the infimum is $0$ and there is no minimum.
